\documentclass[11pt]{article}
\usepackage[a4paper,margin=1in]{geometry}
\usepackage{graphicx}
\usepackage{float}
\usepackage{needspace}
\usepackage{booktabs}
\usepackage{caption}
\usepackage[numbers,sort&compress]{natbib}
\usepackage{hyperref}
\usepackage{amsmath}
\usepackage{longtable}

\title{Adaptation Measures and Cost-Benefit Analysis:\\Brazil Case Study (Porto Alegre)}
\author{
Alisson Tallys Geraldo Fiorentin$^{1}$ \and
Henrique Sanchez Franz$^{2}$ \and
Leonardo Laipelt$^{1}$ \and
Gabriel Matte Rios$^{1}$
}
\date{\today}

\begin{document}
\maketitle

\begin{center}
\footnotesize
$^{1}$Federal University of Rio Grande do Sul (UFRGS) \quad
$^{2}$Federal University of Rio Grande (FURG)
\end{center}
\vspace{1em}

\begin{abstract}
This report documents the extension of the Brazil case study of the IDB RiskHub / CLIMADA
(UNU-EHS) Starter Pack, covering flood risk in Porto Alegre, Rio Grande do Sul, for three
asset types: Schools, Residential, and Companies. Building on the existing Schools example,
the exercise calibrated Residential and Companies, parametrized six adaptation measures per
asset (eighteen in total, covering all six CLIMADA measure mechanisms), and ran a
cost-benefit analysis for the period 2025 to 2050. The cost-benefit analysis uses a
return-period depth hazard, so Residential and Companies are assigned a depth-damage curve
from the literature \citep{huizinga2017}, matched by physical unit to that hazard and
consistent with the curve already used for Schools. Because this curve's amplitude is a
regional literature default, it is calibrated locally by anchoring it to the real loss
observed during the 2024 flood event, using the return-period depth raster whose flooded
extent most closely matches that event as the physical reference point. The resulting
totals are: total climate risk of 320.348 million USD for Schools, 470.267 million USD for
Residential, and 828.925 million USD for Companies. Under this risk estimate, Companies'
six measures remain economically justified, while none of Residential's six measures, at
their currently parametrized cost, pay for themselves. We trace this asymmetry to two
compounding causes: several of Residential's measure costs reflect city-scale programme
budgets that are large relative to its real risk, and Residential's risk, once anchored to
observed loss, is small relative to its very large total exposed value. This is a
decision-relevant finding, not a modelling failure.
\end{abstract}

\newpage
\tableofcontents
\newpage

\section{Introduction}

The IDB RiskHub intersessional assignment asks each participant to work with one or more
calibrated asset types from their country's Starter Pack, select three adaptation measures
relevant to the case study (Brazil participants: three \emph{additional} measures, since
Schools already shipped with three), parametrize those measures in the Entity file with
realistic implementation and maintenance costs and a defined effect on the risk
calculation, run a cost-benefit analysis using the Brazil example notebook, refine the
parametrization until the results appear realistic, and prepare figures suitable for
presentation to decision makers. Participants are further encouraged to calibrate any
remaining uncalibrated assets in their case study, and to note ideas for improving the
Starter Packs.

For Brazil, the three-additional-measures rule applies because the Schools entity already
shipped with three measures (river dredging, levee rebuilding, and an early warning
system). Residential and Companies had no real pre-existing measures, only an inert
placeholder row, so the rule was applied uniformly across all three asset types: three
baseline measures plus three additional measures, six total per asset, eighteen in total.
Applying it uniformly rather than adding three measures only to Schools keeps the three
assets directly comparable, which matters because the most decision-relevant results in this
report come from contrasting how the same six mechanisms perform against three very
different exposure bases.

This report covers four areas of work: the calibration of Residential and Companies, which
was not previously available in the pack; the parametrization of six adaptation measures
per asset, covering all six CLIMADA measure mechanisms; the derivation and local
calibration of the depth-damage curve used for Residential and Companies in the
cost-benefit analysis; and the production of figures and tables suitable for decision
makers. A companion investigation, prompted by the result that none of Residential's
measures pay for themselves, confirmed that this is not a reproducibility issue (Section
\ref{sec:root-cause} and the verification notes in Appendix~\ref{app:notes}).

\section{Methodology}

\subsection{Environment}

The analysis uses CLIMADA version 6.1.0 \citep{climada2026}, pinned deliberately: version 7,
expected later in 2026, rewrites the Measures and CostBenefit modules used throughout this
exercise. Anyone reproducing these results should pin the same version rather than take the
latest release, since the Entity-file and cost-benefit interfaces this exercise depends on
are the ones scheduled to change. The commands below create that environment from scratch
and end with the library's own test suite, which should report \texttt{OK} before any
notebook is run.

\begin{verbatim}
conda create -n climada_env --override-channels -c conda-forge climada=6.1.0 \
    notebook nbconvert -y
conda activate climada_env
python -m unittest climada.engine.test.test_impact   # must end in "OK"
\end{verbatim}

\subsection{Baseline Reproduction}
\label{sec:baseline}

The pack's original notebook, \texttt{09\_costbenefit.ipynb}, was executed unmodified to
confirm the environment. It reproduced the documented reference numbers exactly: total
climate risk 320.348 million USD, average annual risk 19.0853 million USD, residual risk
minus 165.334 million USD, using three measures (river dredging, levee rebuilding, early
warning system) on the Schools asset. As an independent check on this session's own
environment, the same notebook was re-executed a second time, from a pristine clone of the
upstream Starter Pack repository, with zero modifications from any working session. It
reproduced the same three figures exactly, confirming that the CLIMADA library and the
cost-benefit method work correctly and are not sensitive to local environment
peculiarities.

That same check on the upstream repository showed that the Starter Pack, as published,
ships a complete cost-benefit example only for Schools. Residential and Companies exist
upstream only as uncalibrated Entity files, with no measures and no prior cost-benefit
analysis of any kind. The calibration and cost-benefit work for those two assets, described
in the remainder of this report, is therefore a first attempt, not a reproduction of an
existing, previously validated result.

\subsection{Calibration of Residential and Companies}
\label{sec:calibration}

Schools was already calibrated in the pack; Residential and Companies were not. The method
follows the one used in \texttt{06a\_calibration\_schools.ipynb}: a two-dimensional grid
search over the parameters \texttt{max\_damage} and \texttt{max\_intensity}, scaling a
first-guess curve against the 2024 event's flood-duration hazard, followed by a fine
one-dimensional refinement once the coarse grid identifies the neighborhood of the optimum.
The search space is set independently for each asset rather than reused from Schools,
because the three assets differ enormously in exposure scale and in the ratio between
exposed value and observed loss. Residential, for example, is a whole-city grid, with 69.3
billion USD of exposed value against 212 million USD of observed loss, about 0.3 percent,
so the search space is widened to allow \texttt{max\_damage} down to 0.01, well below the
range that would be appropriate for an asset with a higher damage ratio such as Schools.
Companies, with a smaller and more concentrated exposure base, is calibrated over its own
search space in the same way. The resulting fits, summarized in Table \ref{tab:calibration},
reproduce the observed 2024 loss for both assets to within one percent.

\begin{table}[H]
\centering
\caption{Calibration results for Residential and Companies. Only \texttt{max\_damage} is
identified by the data; \texttt{max\_intensity} is not, for the reasons given below.}
\label{tab:calibration}
\small
\begin{tabular}{lrrrr}
\toprule
Asset & max\_damage (identified) & Modelled loss (USD) & Observed loss (USD) & Error \\
\midrule
Residential & 0.028 & 212{,}265{,}415 & 212{,}170{,}680 & 0.04\% \\
Companies & 0.296 & 442{,}034{,}021 & 438{,}755{,}400 & 0.75\% \\
\bottomrule
\end{tabular}
\end{table}

Table \ref{tab:calibration} deliberately reports only \texttt{max\_damage}, even though the
grid search fits two parameters. The reason is that only \texttt{max\_damage} is actually
identified by the data. This is a known property of the method, documented in the pack's
own \texttt{06a} notebook: calibrating two free parameters against a single aggregate
target leaves one degree of freedom, so for every \texttt{max\_intensity} there exists a
\texttt{max\_damage} that reproduces the observation. Our own sweep confirms it directly:
holding \texttt{max\_intensity} fixed at values from 0.25 to 79.75 days, spanning the
entire searched range, a fine search over \texttt{max\_damage} still reaches an error under
3 percent for Residential and under 0.5 percent for Companies at every value tested.

The mechanism behind this degeneracy is specific to the severity of the 2024 event. Median
flood duration at the exposure points is 48 days for Residential and 29 days for Companies,
while the scaled curve reaches its maximum damage at roughly \texttt{max\_intensity}
divided by three. Consequently the great majority of flooded exposed value, about 93
percent for Residential and 91 percent for Companies at the values the search selected,
sits on the flat, saturated part of the curve, where modelled loss no longer depends on
duration at all. The model effectively collapses to the product of \texttt{max\_damage} and
the total flooded exposed value, which pins \texttt{max\_damage} analytically at 0.0277 for
Residential and 0.2942 for Companies, within 1.2 and 0.4 percent of the values in Table
\ref{tab:calibration}, and leaves \texttt{max\_intensity} almost unconstrained.

Two diagnostics confirm that the fitted \texttt{max\_intensity} carries no information. Both
assets settled on opposite edges of the permitted box, Residential at 28.75 days, the
largest value the inherited bound allows, and Companies at 0.25 days, the first point of
the grid: the signature of an optimizer sliding along a flat objective until it meets a
constraint. The Companies value is also physically implausible on its own terms, since a
curve saturating at 0.25 days would imply that a business flooded for two hours suffers the
same damage as one flooded for the full 91 days the event lasted. The stored curves retain
these values (28.75 and 0.25 days) for traceability with the notebooks, but they are
reported here as an artifact of the search, not as an estimate. Resolving the ambiguity
properly would require the same kind of asset-specific expert judgement the pack applies to
Schools, which in turn requires local evidence on how residential and commercial losses
scale with flood duration that this exercise does not have. None of this affects the
cost-benefit results in Section \ref{sec:results}, since the day-based curve is not the one
used there.

These calibrated curves, stored under \texttt{data/entity\_files\_calibrated/}, are valid
for reconstructing the single 2024 event, whose hazard footprint is measured in flood
duration (days). The cost-benefit analysis, by contrast, runs on a return-period depth
hazard (meters); it therefore uses a separate, depth-based damage curve, described in
Section \ref{sec:depth-curve}, matched by physical unit to that hazard. Keeping the two
curves in separate files rather than overwriting one with the other is deliberate: each
remains correct for its own hazard, and the separation makes it impossible to feed a
duration-calibrated curve into a depth-based calculation by accident.

\subsection{Adaptation Measures and Mechanisms}

Each asset received six measures: three that shipped with the pack or follow the same
mechanism family for Residential and Companies, and three additional measures selected for
this exercise, following the intersessional brief's instruction to select measures relevant
to the case study and parametrize their effect on the risk calculation explicitly. Rather
than varying only the numeric parameters of a single mechanism, the eighteen measures were
designed to cover all six adaptation mechanisms CLIMADA supports, spread across the three
assets so that each mechanism is demonstrated at least once. This choice is deliberate: a
decision maker comparing measures that differ only in cost learns less than one comparing
measures that intervene on the risk calculation in genuinely different ways.

\begin{table}[H]
\centering
\caption{CLIMADA measure mechanisms demonstrated in this exercise}
\label{tab:mechanisms}
\begin{tabular}{p{5cm}p{8cm}}
\toprule
Mechanism & Example used \\
\midrule
Linear hazard transform (a, b) & School floor raising, home flood-proofing, portable barriers \\
Hazard frequency cutoff (RP) & Rebuilding levees, city-wide drainage upgrade, commercial levee \\
Risk transfer (insurance) & Portfolio, residential parametric, and business interruption insurance \\
Exposures swap (relocation) & Partial relocation in all three assets \\
Linear MDD transform & Early warning system, community warning, elevated inventory storage \\
Linear PAA transform & Flood-resistant building codes, business continuity planning \\
\bottomrule
\end{tabular}
\end{table}

For insurance, attachment and cover limits were derived from the model's own exceedance
curve (\texttt{Impact.calc\_freq\_curve}, using the 10-year return period as attachment and
the 100-year return period as cover), not from arbitrary values. The cost loading applied
was 1.5 times the pure premium, which is why the insurance benefit-to-cost ratio sits close
to 1 divided by 1.5, approximately 0.66, across all three assets. This value reflects the
cost of transferring volatility, not an expectation of savings.

For relocation, an HDF5 Exposures file was built in which the locations with the highest
depth at the 100-year return period have their value reduced to 40 percent of the original,
corresponding to a partial relocation of 60 percent of the exposed value. A second variant
was built for 2050, already scaled by the 2 percent annual economic growth rate; without
this adjustment, the measure's benefit would incorrectly include the removal of economic
growth itself from the future scenario.

\subsection{Depth-Damage Curve for the Cost-Benefit Analysis}
\label{sec:depth-curve}

Residential and Companies are assigned a depth-based damage curve for the cost-benefit
analysis, matched to the return-period depth hazard used there. Its shape comes from the
curve published by the Joint Research Centre \citep{huizinga2017}: Table 3-5, for
residential use in South America, is adopted for Residential, and Table 3-9, for commercial
use in South America, is adopted for Companies. This is the same literature source and the
same pattern already used, though without explicit citation, by the Schools entity shipped
with the pack, which uses an equivalent depth-based curve.

Because a regional literature curve's amplitude is not, by itself, validated for Porto
Alegre, it is calibrated locally against the only available local data point: the real loss
observed during the 2024 flood event. No depth-based footprint exists for that event, only
a duration-based one, so a physical proxy is needed. Among the available return-period
depth rasters, the one whose flooded extent most closely matches the extent actually
observed in 2024 serves as that proxy. Comparing flooded area across rasters (pixel count
with depth greater than zero, converted to square kilometers accounting for Porto Alegre's
latitude), the 2024 event's extent (141.2 km\textsuperscript{2}) falls almost exactly
between the 100-year return period (134.3 km\textsuperscript{2}) and the 500-year return
period (148.0 km\textsuperscript{2}), closest to the 200-year return period (140.9
km\textsuperscript{2}, a difference of only 0.3 km\textsuperscript{2}). This is consistent
with public accounts of the 2024 Rio Grande do Sul floods as an exceptional, near
record-breaking event for Porto Alegre, and the 200-year raster is used as the anchor.

The curve's amplitude is calibrated by rescaling: keeping the same depth breakpoints and
the same relative shape of the JRC curve, every mean damage degree (MDD) value is
multiplied by a single scale factor, computed as the ratio of the observed 2024 loss to the
loss the unscaled curve would produce on the 200-year return-period raster (0.048339 for
Residential, 0.360351 for Companies). Because no value of the rescaled curve saturates at
MDD equal to one within the observed depth range (0 to about 9 meters), this scaling
propagates exactly linearly through the rest of the analysis: total risk, average annual
risk, and each measure's benefit (with the exception of insurance) are each the value the
unscaled curve would have produced, multiplied by the same asset-specific scale factor.
Insurance attachment and cover limits are computed from this calibrated curve's own
exceedance curve, following the same method as Section 2.4.

The resulting totals are 470.267 million USD for Residential and 828.925 million USD for
Companies (Section \ref{sec:results}). Measure costs are parametrized independently, as
construction or municipal programme budgets, not as a quantity derived from the risk scale;
Section \ref{sec:root-cause} examines the consequence of that independence in detail: none
of Residential's six measures pay for themselves at their parametrized cost, while all six
of Companies' measures continue to. This is an explicit limitation to keep in mind: the
curve is anchored to a single aggregate data point (one event's total loss), not to a
spatially distributed calibration like the one performed for the duration hazard in Section \ref{sec:calibration}, the same caveat that applies to any calibration built on few reference points.

As an independent check that this scale factor is not an arbitrary or forced choice, the
same calibration was reproduced using CLIMADA's own calibration framework
(\texttt{climada.util.calibrate}), which minimizes the mean squared error between modelled
and observed loss over a free scale parameter using \texttt{scipy.optimize}. Three
optimization algorithms (Nelder-Mead, L-BFGS-B, and Powell) were run independently and all
converged to the same values obtained analytically, 0.048339 for Residential and 0.360351
for Companies, confirming that these are the cost-minimizing values for this one-parameter
model given the observed 2024 loss, not a value chosen to produce a particular outcome.

\subsection{Sensitivity to the Anchor Return Period}
\label{sec:sensitivity}

The scale factor is sensitive to which return period is used as the physical proxy for the
2024 event: computed the same way at every available return period, it ranges from 2.37
(Residential, RP10) down to 0.039 (Residential, RP500), since the raw curve's modelled loss
grows with depth while the observed 2024 loss is fixed. This makes the RP200 choice, based
on flooded-area similarity (Section \ref{sec:depth-curve}), the single most consequential
methodological decision in this analysis, and it is worth checking how much the reported
results would change under a different, still physically plausible choice.

\begin{table}[H]
\centering
\caption{Total risk and non-insurance measure B/C under alternative anchor return periods}
\label{tab:sensitivity}
\small
\begin{tabular}{lrrr}
\toprule
& RP100 & RP200 (used) & RP500 \\
\midrule
Residential total risk (M USD) & 569.0 & 470.3 & 378.1 \\
Residential measures, B/C range & 0.23 to 0.88 & 0.19 to 0.73 & 0.15 to 0.59 \\
Companies total risk (M USD) & 1{,}019.8 & 828.9 & 632.5 \\
Companies measures, B/C range (excl. relocation) & 4.37 to 11.17 & 3.55 to 9.08 & 2.71 to 6.93 \\
Companies relocation, B/C & 1.62 & 1.32 & 1.01 \\
\bottomrule
\end{tabular}
\end{table}

Across this entire range, spanning a 100-year to 500-year event, the decision-relevant
conclusion is stable: every one of Residential's non-insurance measures stays below a
benefit-to-cost ratio of 1, and every one of Companies' measures other than relocation
stays above 1. Companies' relocation measure is the one result sensitive enough to flag on
its own: its ratio moves from 1.62 at RP100 to 1.01 at RP500, staying on the right side of
break-even across the tested range but by a narrow and anchor-dependent margin. Insurance's
ratio is excluded from this table because it depends on the model's exceedance curve rather
than scaling linearly with the anchor choice; in practice it stays close to 0.66 under every
return period tested, since it is set by the 1.5 times cost loading rather than by the
underlying risk level.

\subsection{Acceptance Criteria}
\label{sec:refinement}

Each measure's parametrization is checked against three criteria: no infinite, undefined,
or greater-than-50 benefit-to-cost ratio, which would signal an unrealistically low cost;
no individual measure's benefit exceeding the asset's total climate risk; and insurance
attachment and cover limits always derived from the model, never from arbitrary estimates.
The first two are sanity checks against a mis-specified cost or a benefit computed against
the wrong baseline; the third exists because risk-transfer parameters chosen by hand can
produce almost any premium, and with it almost any benefit-to-cost ratio, without the error
being visible anywhere in the output.

Under the final parametrization, no benefit-to-cost ratio exceeds 9.08, well below the
ceiling of 50, and no individual benefit exceeds the asset's total climate risk; both
sanity criteria are satisfied. For Residential, all six ratios fall below 1.0, which is not
a sign of a parametrization error: it is the honest conclusion that, at this risk estimate,
these specific measures, at the costs they were parametrized with, do not pay for
themselves within the study horizon. Adjusting the costs to force ratios above one would
substitute the desired outcome for the best available estimate of cost and risk, which is
exactly what these criteria are meant to prevent.

\section{Results}
\label{sec:results}

Each of the three subsections below covers one asset type in turn: Schools, Residential,
and Companies. Every subsection follows the same structure, so results across assets can
be compared directly. It opens with the asset's total climate risk and average annual
risk for 2050, then a table listing all six adaptation measures with their mechanism,
cost, benefit, and benefit-to-cost ratio, then the five figures generated for that asset
(the cost-benefit chart, benefits by return period, the waterfall of averted risk, the
expected annual impact map, and the depth map with relocation points), and closes with a
paragraph interpreting what those figures show. A final subsection, Cross-Asset Synthesis,
draws out the patterns that only become visible when comparing all three assets side by
side, such as the stability of the insurance ratio and the climate-versus-growth split.

\subsection{Schools}

Schools is the one asset in this exercise whose calibration and damage curve were never
touched: its risk figures below are exactly those the Starter Pack shipped with, unaffected
by the calibration and curve-derivation work described in Section 2 for Residential and
Companies. Total climate risk is 320.348 million USD, with an average annual risk of
19.0853 million USD; both numbers are reproduced exactly by the pristine, unmodified
upstream notebook, as described in Section \ref{sec:baseline}. Table \ref{tab:schools}
lists all six measures, three that shipped with the pack (river dredging, levee
rebuilding, and the early warning system) and three added for this exercise (school floor
raising, portfolio insurance, and partial relocation), spanning five of the six CLIMADA
measure mechanisms on their own.

\begin{table}[H]
\centering
\caption{Schools: adaptation measures and cost-benefit results}
\label{tab:schools}
\begin{tabular}{p{3.6cm}p{3cm}rrr}
\toprule
Measure & Mechanism & Cost (M USD) & Benefit (M USD) & B/C \\
\midrule
River dredging & hazard $-1.5$ m & 789.562 & 195.782 & 0.25 \\
Rebuilding levees & RP cutoff $\approx$ 100y & 161.171 & 257.866 & 1.60 \\
Early warning system & MDD $\times 0.9$ & 1.500 & 32.035 & 21.36 \\
School floor raising & hazard $-0.4$ m & 7.625 & 64.244 & 8.43 \\
Portfolio flood insurance & risk transfer 32.3M xs 116.9M & 258.463 & 171.475 & 0.66 \\
Partial relocation (5 of 122) & exposures swap & 10.500 & 34.764 & 3.31 \\
\bottomrule
\end{tabular}
\end{table}

Costs in Table \ref{tab:schools} range from 1.5 million USD for the early warning system to
789.6 million for river dredging, and benefit-to-cost ratios range even more widely, from
0.25 for dredging up to 21.36 for the early warning system, the two ends of a pattern that
recurs across all three assets: hazard-transform measures with no protection ceiling are
expensive relative to the risk they avert, while cheap, narrowly targeted measures post the
highest ratios precisely because their cost base is so small. The five figures that follow
show this pattern from different angles, from the raw cost-benefit comparison to how each
measure's contribution changes with event severity.

\needspace{3.08in}
\begin{figure}[H]
\centering
\includegraphics[width=0.55\textwidth]{figures/brazil_schools_cost_benefit.png}
\caption{Cost-benefit results, Schools}
\label{fig:schools-cb}
\end{figure}

Figure \ref{fig:schools-cb} plots each measure as a rectangle whose width is the net present
value of averted damage and whose height is the benefit-to-cost ratio, so area is
proportional to neither quantity alone and the two dimensions must be read together. The
early warning system is the tall, narrow bar on the left: an excellent ratio built on a very
small absolute benefit. River dredging is the opposite, the widest bar in the chart but so
flat that it sits below the break-even line. For a decision maker this is the central
trade-off in the Schools portfolio, since the measure that averts the most damage in
absolute terms is also the one that fails to justify its cost.

\needspace{2.73in}
\begin{figure}[H]
\centering
\includegraphics[width=0.55\textwidth]{figures/brazil_schools_benefits_by_return_period.png}
\caption{Benefits by return period, Schools}
\label{fig:schools-rp}
\end{figure}

Figure \ref{fig:schools-rp} breaks the same benefits down by event severity, and explains
why dredging looks so different from the rest. As a hazard transform with no protection
ceiling, it keeps growing in absolute terms as events get more severe, and it dominates the
benefit at the 200-year return period. Levees, with a cutoff near the 100-year return
period, and insurance between them absorb most of the average risk across the middle of the
range. Relocation and the early warning system appear only as thin bands throughout, which
is consistent with their small absolute benefits even where their ratios are favourable.

\needspace{3.29in}
\begin{figure}[H]
\centering
\includegraphics[width=0.48\textwidth]{figures/brazil_schools_waterfall_averted_risk.png}
\caption{Waterfall of averted risk, Schools}
\label{fig:schools-wf}
\end{figure}

Figure \ref{fig:schools-wf} separates the drivers of future risk rather than the effect of
the measures. Of the 19 million USD of risk projected for 2050, 6 million already existed in
2025, a further 4 million comes from economic growth, and a further 10 million comes from
climate change. Roughly 71 percent of the growth in risk between now and 2050 is therefore
attributable to the climate scenario rather than to more assets being built in exposed
places, a proportion that recurs closely in the other two assets and is discussed in Section
\ref{sec:synthesis}.

\needspace{4.26in}
\begin{figure}[H]
\centering
\includegraphics[width=0.55\textwidth]{figures/brazil_schools_map_expected_annual_impact.png}
\caption{Expected annual impact map, Schools}
\label{fig:schools-eai}
\end{figure}

Figure \ref{fig:schools-eai} moves from aggregate to spatial, showing where the expected
annual impact actually falls. Risk concentrates in the Farrapos and Humaita neighborhoods
and around Sarandi, the lowest-lying fringes along Lake Guaiba, with the rest of the school
network contributing very little. Because Schools is a point-facility dataset of 122
locations rather than a grid, individual high-risk schools are visible as discrete markers,
which is what makes a targeted measure such as relocation feasible to specify in the first
place.

\needspace{4.99in}
\begin{figure}[H]
\centering
\includegraphics[width=0.55\textwidth]{figures/brazil_map_rp100_depth_schools.png}
\caption{100-year return period depth and relocation points, Schools}
\label{fig:schools-depth}
\end{figure}

Figure \ref{fig:schools-depth} overlays the 100-year flood depth on the school locations and
marks the five schools selected for partial relocation. They sit exactly where the depth
footprint is darkest, which is the criterion by which they were chosen, with one isolated
case further south along the lake. The clustering explains why relocating only five of 122
schools still returns a benefit-to-cost ratio above three: risk in this asset is
concentrated enough that a small, well-targeted intervention captures a disproportionate
share of it.

\subsection{Residential}

Residential is the asset most affected by the depth-damage curve derivation in Section
\ref{sec:depth-curve}, since it is calibrated and risk-assessed for the first time in this
exercise rather than reusing an existing pack example. Total climate risk is 470.267
million USD, with an average annual risk in 2050 of 28.0255 million USD, figures already
anchored to Porto Alegre's real 2024 loss and checked for sensitivity to the anchor choice
in Section \ref{sec:sensitivity}. Table \ref{tab:residential} lists the six measures
parametrized for this asset: home flood-proofing and a city-wide drainage upgrade, both
carried over in mechanism from the Schools example; a parametric insurance layer sized to
this asset's own exceedance curve; partial relocation of the highest-risk households; and
two measures added specifically for Residential, a community-based early warning system
and flood-resistant building codes for renovations.

\begin{table}[H]
\centering
\caption{Residential: adaptation measures and cost-benefit results}
\label{tab:residential}
\begin{tabular}{p{3.6cm}p{3cm}rrr}
\toprule
Measure & Mechanism & Cost (M USD) & Benefit (M USD) & B/C \\
\midrule
Home flood-proofing & hazard $-0.3$ m & 362.064 & 70.321 & 0.19 \\
Residential parametric insurance & risk transfer 47.8M xs 175.4M & 382.446 & 251.631 & 0.66 \\
City-wide drainage upgrade & RP cutoff $=20$y & 300.000 & 218.109 & 0.73 \\
Relocation (50 of 20{,}187) & exposures swap & 8.250 & 3.924 & 0.48 \\
Community-based early warning & MDD $\times 0.85$ & 150.000 & 70.540 & 0.47 \\
Flood-resistant building codes & PAA $\times 0.8$ & 200.000 & 94.053 & 0.47 \\
\bottomrule
\end{tabular}
\end{table}

None of these six measures pays for itself at its parametrized cost, at this risk estimate.
This is the central decision-relevant finding of this report, and it is examined in detail
in Section \ref{sec:root-cause}, which traces it to two compounding causes: the scale of
several of these measures' budgets relative to Residential's real risk, and the underlying
gap between how much value is exposed and how much of that value is actually at risk once
the damage curve is anchored to a real observed loss rather than left at its literature
amplitude. The five figures below show the same conclusion from different angles, from the
raw cost-benefit comparison to how the underlying risk itself is distributed across the
city.

\needspace{3.10in}
\begin{figure}[H]
\centering
\includegraphics[width=0.55\textwidth]{figures/brazil_residential_cost_benefit.png}
\caption{Cost-benefit results, Residential}
\label{fig:residential-cb}
\end{figure}

Figure \ref{fig:residential-cb} makes the headline result immediately visible: every
rectangle sits below the break-even line, so no measure in this portfolio returns more than
it costs. The chart is worth reading alongside the Schools equivalent (Figure
\ref{fig:schools-cb}), where the bars span a wide range of heights. Here they are compressed
into a narrow band between 0.19 and 0.73, which is the visual signature of measure budgets
that are large relative to the risk they address rather than of any single mis-specified
measure.

\needspace{2.66in}
\begin{figure}[H]
\centering
\includegraphics[width=0.55\textwidth]{figures/brazil_residential_benefits_by_return_period.png}
\caption{Benefits by return period, Residential}
\label{fig:residential-rp}
\end{figure}

Figure \ref{fig:residential-rp} shows a sharp cutoff pattern. The city-wide drainage
upgrade, with its cutoff at the 20-year return period, covers the entirety of risk at the 5
and 10-year return periods, below its protection ceiling, then disappears completely from
the 20-year return period onward once events exceed that ceiling. Beyond that point
insurance, sitting at the highest risk layer, and the building codes take over. At the
200-year return period a slice remains uncovered between the top of the stacked measures and
the total risk envelope, showing that the six measures together still leave a visible
residual risk against the most extreme event.

\needspace{3.32in}
\begin{figure}[H]
\centering
\includegraphics[width=0.48\textwidth]{figures/brazil_residential_waterfall_averted_risk.png}
\caption{Waterfall of averted risk, Residential}
\label{fig:residential-wf}
\end{figure}

Figure \ref{fig:residential-wf} decomposes the growth in risk to 2050. Risk rises from 8.2
million USD in 2025 to 28.0 million in 2050, an increase of 5.3 million attributable to
economic growth and 14.5 million to climate change, about 73 percent climate against 27
percent growth. That this matches the Schools split so closely, despite a completely
different exposure base and calibration, is one of the more robust findings of the exercise
and is taken up in Section \ref{sec:synthesis}.

\needspace{4.20in}
\begin{figure}[H]
\centering
\includegraphics[width=0.55\textwidth]{figures/brazil_residential_map_expected_annual_impact.png}
\caption{Expected annual impact map, Residential}
\label{fig:residential-eai}
\end{figure}

Figure \ref{fig:residential-eai} shows a spatial pattern quite distinct from Schools.
Instead of a handful of concentrated points, residential exposure covers nearly the entire
urban footprint of Porto Alegre, because this asset is a city-wide gridded dataset of more
than 20,000 points rather than a list of facilities. This breadth is precisely what makes
city-scale measures expensive here: an intervention defined across the whole grid has to be
funded across the whole grid, while the risk it addresses remains concentrated in the
low-lying north.

\needspace{4.99in}
\begin{figure}[H]
\centering
\includegraphics[width=0.55\textwidth]{figures/brazil_residential_map_rp100_depth_relocation.png}
\caption{100-year return period depth and relocation points, Residential}
\label{fig:residential-depth}
\end{figure}

Figure \ref{fig:residential-depth} marks the 50 grid points selected for partial relocation
against the 100-year depth footprint. Unlike Schools, where the selected sites form a tight
cluster, these spread across several riverside neighborhoods with additional pockets further
south. The wider spread is one reason relocation performs worse here than for the other two
assets: the same conservative rule, moving 60 percent of value at the deepest points,
captures a smaller share of total risk when that risk is diffuse rather than concentrated.

\subsection{Companies}

Companies shares its calibration method and depth-damage curve derivation with Residential,
but differs sharply in how concentrated its exposure is: a smaller, denser commercial base
along the riverfront and in the historic center, rather than a grid spanning the whole city.
Total climate risk is 828.925 million USD, with an average annual risk in 2050 of 49.6197
million USD, both already anchored to the observed 2024 loss in the same way as Residential.
Table \ref{tab:companies} lists the six measures parametrized for this asset: portable flood
barriers and a levee protecting the commercial district, mirroring Residential's
hazard-transform and cutoff mechanisms; business interruption insurance; relocation of the
highest-risk businesses; elevated inventory storage; and business continuity and supply
chain diversification planning.

\begin{table}[H]
\centering
\caption{Companies: adaptation measures and cost-benefit results}
\label{tab:companies}
\begin{tabular}{p{3.6cm}p{3cm}rrr}
\toprule
Measure & Mechanism & Cost (M USD) & Benefit (M USD) & B/C \\
\midrule
Portable flood barriers & hazard $-0.3$ m & 36.533 & 147.182 & 4.03 \\
Business interruption insurance & risk transfer 66.7M xs 356.6M & 713.520 & 474.347 & 0.66 \\
Relocation (20 of 4{,}366) & exposures swap & 10.250 & 13.508 & 1.32 \\
Levee (commercial district) & RP cutoff $=20$y & 50.000 & 321.952 & 6.44 \\
Elevated inventory storage & MDD $\times 0.8$ & 18.266 & 165.785 & 9.08 \\
Business continuity and supply chain diversification & PAA $\times 0.7$ & 70.000 & 248.677 & 3.55 \\
\bottomrule
\end{tabular}
\end{table}

Unlike Residential, all six measures pay for themselves (B/C $>1$), a direct consequence of
Companies' much smaller, more concentrated exposure base relative to its risk (Section
\ref{sec:root-cause} quantifies this contrast directly). Relocation, at 1.32, does so with
the narrowest margin of the six, and is the one result in this report sensitive enough to
the anchor return period (Section \ref{sec:sensitivity}) to be worth flagging as less
certain than the rest; the other five measures stay comfortably above break-even across the
full range of anchor choices tested. The five figures below show this pattern in more
detail, from the raw cost-benefit comparison through to how the risk itself is distributed
across the city.

\needspace{3.11in}
\begin{figure}[H]
\centering
\includegraphics[width=0.55\textwidth]{figures/brazil_companies_cost_benefit.png}
\caption{Cost-benefit results, Companies}
\label{fig:companies-cb}
\end{figure}

Figure \ref{fig:companies-cb} is the visual counterpart to Figure \ref{fig:residential-cb}
and the contrast is stark: here the rectangles rise well above the break-even line, with
elevated inventory storage reaching 9.08 and the commercial-district levee 6.44. The same
six mechanisms, applied to a much smaller and more concentrated exposure base, produce a
portfolio that pays for itself several times over. Only insurance sits below the line, and
for the structural reason discussed in Section \ref{sec:synthesis} rather than because of
how it was costed.

\needspace{2.63in}
\begin{figure}[H]
\centering
\includegraphics[width=0.55\textwidth]{figures/brazil_companies_benefits_by_return_period.png}
\caption{Benefits by return period, Companies}
\label{fig:companies-rp}
\end{figure}

Figure \ref{fig:companies-rp} repeats the cutoff pattern seen for Residential. The levee,
with a cutoff at the 20-year return period, dominates the benefit at the 5 and 10-year
return periods and drops to zero from 20 years onward. Beyond that point insurance dominates
at the highest risk layer, while the MDD and PAA transform mechanisms, elevated inventory
storage and business continuity planning, contribute a relatively stable middle band across
every return period, since neither has a protection ceiling that the event can exceed.

\needspace{3.32in}
\begin{figure}[H]
\centering
\includegraphics[width=0.48\textwidth]{figures/brazil_companies_waterfall_averted_risk.png}
\caption{Waterfall of averted risk, Companies}
\label{fig:companies-wf}
\end{figure}

Figure \ref{fig:companies-wf} completes the set of three risk decompositions. Risk rises
from 14.3 million USD in 2025 to 49.6 million in 2050, with 9.1 million attributable to
economic growth and 26.2 million to climate change, close to 74 percent climate against 26
percent growth. This is the third asset to land in the same narrow band despite entirely
independent exposure data, calibration, and measures, which is the strongest internal
consistency check available in this exercise.

\needspace{4.19in}
\begin{figure}[H]
\centering
\includegraphics[width=0.55\textwidth]{figures/brazil_companies_map_expected_annual_impact.png}
\caption{Expected annual impact map, Companies}
\label{fig:companies-eai}
\end{figure}

Figure \ref{fig:companies-eai} shows a double concentration rather than the single cluster
of Schools or the diffuse spread of Residential. A dense band runs along the river through
the Farrapos, Humaita, and Sarandi neighborhoods, coinciding with the lowest and most
flood-prone ground, and a second, much smaller cluster sits in Porto Alegre's historic
center, its commercial and port district. That commercial activity concentrates in exactly
the places most exposed to flooding is what gives this asset its high risk-to-exposure
ratio, and in turn why targeted measures perform so well against it.

\needspace{4.99in}
\begin{figure}[H]
\centering
\includegraphics[width=0.55\textwidth]{figures/brazil_companies_map_rp100_depth_relocation.png}
\caption{100-year return period depth and relocation points, Companies}
\label{fig:companies-depth}
\end{figure}

Figure \ref{fig:companies-depth} marks the 20 businesses selected for relocation against the
100-year depth footprint. Almost all fall inside the dense riverside band, with only a few
isolated points outside it, an intermediate case between the tight clustering of Schools and
the wide dispersion of Residential. This is consistent with relocation returning a positive
but narrow benefit-to-cost ratio of 1.32 here, the one Companies result flagged in Section
\ref{sec:sensitivity} as sensitive to the anchor choice.

\subsection{Cross-Asset Synthesis}
\label{sec:synthesis}

Insurance stays consistently close to a benefit-to-cost ratio of 0.66 across all three
assets. This is a mechanical consequence of the uniformly applied 1.5 times cost loading,
not a modelling failure: insurance's value lies in transferring volatility, not in expected
long-run savings, and this ratio depends only on the loading factor, not on the risk scale.
It is worth stating explicitly to a decision maker, because a bar chart that places
insurance below every other measure invites the reading that it was parametrized badly, when
in fact a ratio near the reciprocal of the loading is exactly what a correctly specified
risk-transfer layer should produce.

Wide-infrastructure measures, such as levees, dredging, and drainage, show very different
benefit-to-cost ratios depending on the asset evaluated. In Schools, these measures protect
the entire city but are evaluated only against the narrow slice of school assets, yielding
ratios between 0.25 and 1.60. In Companies, the commercial-district levee, evaluated against
the population it actually protects, reaches 6.44. In Residential, the same "broader base"
logic applies to the drainage upgrade, but its risk scale leaves no margin: its ratio falls
to 0.73. In other words, evaluating against a broader base continues to raise the
benefit-to-cost ratio relative to Schools, but is not, by itself, enough to offset a real
risk much smaller than the exposed value it is measured against.

Relocation is not uniformly positive across the three assets: in Schools the ratio is 3.31;
in Companies it is 1.32, positive but by a narrow margin; in Residential it is 0.48, not
paying for itself. The conservative approach adopted, relocating only 60 percent of the
exposed value at the highest-depth points, is still not enough to justify Residential's
relocation cost against its estimated risk. The pattern across the three assets tracks how
spatially concentrated each one's risk is, as the depth maps in Section \ref{sec:results}
show: tightly clustered for Schools, intermediate for Companies, and diffuse for
Residential. Relocation rewards concentration, because a fixed number of relocated points
captures a larger share of total risk when that risk is not spread thinly.

The proportion between climate change and economic growth in future risk growth is notably
stable across all three assets, between 71 and 75 percent attributable to climate change and
between 25 and 29 percent to economic growth. This is a sign of the exercise's robustness:
the three assets use completely different exposure bases, calibrations, and measures, but
share the same future climate scenario, the 2025 to 2050 return-period remapping, and the
same 2 percent annual economic growth rate. The stability of this proportion suggests that
the climate scenario, not particular parametrization choices, is the main driver of
projected future risk growth.

Finally, measures with a return-period cutoff, such as drainage and levees, always drop to
zero above their protection ceiling. This pattern appears clearly in the
benefits-by-return-period charts for Residential and Companies: full protection up to the
20-year return period, zero contribution afterward. This is a known limitation of that
mechanism, offering complete protection or none, with no partial gradation, and it is worth
explaining clearly to a decision maker before the measure is adopted in isolation.

\section{Root Cause Analysis: Why Residential Measures Do Not Pay Off}
\label{sec:root-cause}

The result that none of Residential's six measures pay for themselves, while all six of
Companies' do, prompted a dedicated investigation, since a result that surprising warrants
verification before it is trusted. Two checks were performed. First, a pristine clone of
the upstream Starter Pack repository confirmed that no official Residential or Companies
cost-benefit example exists anywhere upstream: only Schools ships a complete, calibrated
example. There is therefore no previously "working" Residential result that this exercise
could have regressed; this is the first time this calculation has ever been performed for
that asset. Second, the notebook re-execution described in Section \ref{sec:baseline} confirmed the
pipeline itself reproduces its own numbers exactly, ruling out an execution or environment
bug.

With reproducibility and baseline validity confirmed, the asymmetry itself still needs an
explanation, since it is large enough that a reader could reasonably suspect a modelling
error rather than a genuine finding. The analysis below identifies two compounding causes,
one about how the six measures were costed and one about the underlying economics of the
two exposure bases, and quantifies each rather than leaving the explanation qualitative.
Table \ref{tab:rootcause} sets out the comparison that motivates both: Residential and
Companies differ enormously in exposed value, but far less in total climate risk, and that
gap alone accounts for most of what follows.

\begin{table}[H]
\centering
\caption{Cost-to-risk comparison, Residential versus Companies}
\label{tab:rootcause}
\small
\begin{tabular}{lp{2cm}p{2cm}p{2cm}p{2.2cm}p{1.5cm}}
\toprule
Asset & Exposed value (2025) & Total climate risk & Risk / exposed value & Sum of 6 measure costs & Cost / risk \\
\midrule
Residential & 69.312 bn USD & 470.267 M USD & 0.68\% & 1{,}402.760 M USD & 298\% \\
Companies & 3.656 bn USD & 828.925 M USD & 22.67\% & 898.569 M USD & 108\% \\
\bottomrule
\end{tabular}
\end{table}

\textbf{Cause 1: several of Residential's measure costs are large, programme-scale budgets
relative to its risk.} Three of Residential's six measures, Home flood-proofing (362.064
million USD), Community-based early warning (150.000 million USD), and Flood-resistant
building codes (200.000 million USD), are costed as city-scale programmes intended to reach
the full residential base, 69.3 billion USD of exposed value spread across more than 20,000
grid points. Only one of Companies' six measures is costed at a comparable programme scale.
This asymmetry in cost scale, independent of any curve calibration question, already
accounts for much of the difference in benefit-to-cost outcome between the two assets.

\textbf{Cause 2: the scale mismatch is also structural.} Residential has about 19 times
more exposed value than Companies, but its total climate risk is only marginally higher in
absolute dollar terms. This makes Residential's risk intensity, risk as a share of exposed
value, about 33 times smaller than Companies' (0.68 percent against 22.67 percent). The sum
of Residential's six measure costs, 1.4 billion USD, almost three times its own total risk,
simply does not fit inside the risk it addresses, even where no individual cost is
unusually large in absolute terms.

The difference in how much the unscaled literature curve overestimated the observed 2024
loss, roughly 20.7 times for Residential against 2.8 times for Companies (Section
\ref{sec:depth-curve}), traces back to the JRC curve itself: the "South America
residential" table fit Porto Alegre's real 2024 loss far worse than the "South America
commerce" table did. It is not possible to say with confidence, from the data currently
available, whether this is because the curve is a poor match for Porto Alegre's residential
building stock specifically, or because that stock is itself atypical of the regional
average the curve represents; resolving this would require a depth-based footprint of the
2024 event, the same limitation already noted in Section \ref{sec:depth-curve}.

This is not a computational error. It is a legitimate, if uncomfortable, finding: at the
best available estimate of real flood risk, Residential adaptation as currently scoped and
costed in this exercise is not economically justified, while Companies' portfolio remains
justified with reduced margins. The practical implication is not that residential adaptation
should be abandoned, but that it should be re-scoped: the measures here were specified at
city scale, and the analysis suggests the case for them has to be made either at a smaller,
better-targeted scale or on grounds other than avoided direct economic damage, such as
displacement, health, or continuity of services, none of which this model prices.

\section{Discussion and Conclusions}

The exercise confirms that extending a cost-benefit example from one asset to others within
the same case study is not a mechanical exercise in copying code. Calibrating a new asset
requires a search space of its own, sensitive to that asset's exposure scale, and
integrating calibration with cost-benefit analysis requires explicit attention to the
physical unit of the hazard used at each step: Residential and Companies are calibrated
against a duration hazard for event reconstruction, but the cost-benefit analysis runs on a
depth hazard, so it uses a separate, depth-based curve (Section \ref{sec:depth-curve}). A
declared physical unit compatible with the hazard actually used is necessary but not
sufficient: the literature curve's amplitude, tested with no adjustment against the real
loss observed in 2024, overestimated that loss by 20.7 times for Residential and 2.8 times
for Companies, which is why its amplitude is calibrated locally rather than taken at face
value. The broader lesson is that a damage curve's physical unit and its magnitude are two
separate properties to verify; a curve can be dimensionally correct and still produce
substantially incorrect results, silently, because the output looks plausible enough not to
raise suspicion unless it is deliberately checked against a real observed loss.

From a decision-making standpoint, the final picture is more heterogeneous than a single
conclusion of the form "adaptation measures are worthwhile." For Companies, the six-measure
portfolio makes economic sense, with modest margins but all above break-even. For
Residential, none of the six parametrized measures pays for
itself against the real estimated risk: not the wide-infrastructure measure (drainage), not
the partial relocation, not the MDD or PAA measures. This does not mean adaptation is not
worthwhile for Porto Alegre's residential sector; it means the measures \emph{as
parametrized in this exercise}, at these costs and this scale, are not justified against the
best available estimate of risk. A future refinement would need to reconsider the scale of
the interventions, for example targeting only the highest-risk points rather than the whole
city, or seek genuine cost reductions, rather than simply accepting the current figures.
Parametric insurance, across all three assets, remains structurally the lowest
benefit-to-cost mechanism (approximately 0.66), a cost of transferring volatility worth
explaining to a decision maker as expected, not as a poorly parametrized measure.

The most important limitation of this exercise is that the final damage curve, though
anchored to the real 2024 loss, was calibrated against a single aggregate data point, the
event's total loss, not against a spatial distribution of losses like the one that guided
the duration hazard's calibration in Section \ref{sec:calibration}. The 200-year return period used as a
proxy for 2024's depth was chosen by area similarity between two distinct hazard models,
the 2024 duration model and the return-period depth model, not by a direct physical
correspondence; it is the best available approximation, not a calibration with the same
rigor as Section \ref{sec:calibration}. Any future refinement of the Brazil case study should prioritize
obtaining or reconstructing a depth-based footprint of the 2024 event, to enable a direct,
spatially distributed calibration, the way one already exists for the duration hazard.

\section{Recommendations for Replication}

Every number in this report is reproducible from the data and code delivered alongside it,
without needing to rebuild anything from scratch. The five notebooks that produce the
results in Section \ref{sec:results}, the calibration in \texttt{06b} and the four
cost-benefit runs in \texttt{09} through \texttt{12}, can each be re-executed independently
once the environment described in Section 2.1 is set up, and each writes its own figures
and printed totals in place, so a discrepancy in any one number can be traced back to the
specific notebook cell that produced it. The commands below assume a terminal open in the
root of this repository, with the \texttt{climada\_env} environment already created and
activated.

\begin{verbatim}
conda activate climada_env
cd notebooks/
jupyter nbconvert --to notebook --execute --inplace \
    06b_calibration_residential_companies.ipynb
jupyter nbconvert --to notebook --execute --inplace 09_costbenefit.ipynb
jupyter nbconvert --to notebook --execute --inplace \
    10_costbenefit_additional_measures.ipynb
jupyter nbconvert --to notebook --execute --inplace \
    11_costbenefit_residential.ipynb
jupyter nbconvert --to notebook --execute --inplace \
    12_costbenefit_companies.ipynb
\end{verbatim}

All required input files, including the calibrated and measure-parametrized Entity files
and the relocation HDF5 files, are already available in \texttt{data/}. It is not necessary
to run the numbered \texttt{.py} scripts to reproduce the results in this report; they
exist only as a record of how each file was built and are documented in
\texttt{notebooks/README.md}. Running them is only necessary to rebuild the inputs from
scratch or to adapt part of the exercise.

To reproduce or extend this analysis within the Brazil case study, five general steps
apply. First, calibrate the asset of interest following the method of notebooks
\texttt{06a} and \texttt{06b}, adjusting the search space to that asset's own exposure
scale, since the original search space does not transfer between assets of different
scales. Second, and most important, confirm that the physical unit of the hazard used in
calibration matches the physical unit of the hazard that will be consumed in the
cost-benefit analysis; this should be a mandatory step before any calibration, since a unit
mismatch produces no exceptions and no error messages, only numerically plausible and
substantially incorrect results. Third, if a literature curve is substituted in to resolve
a unit or domain mismatch, validate its magnitude against the best available local loss
data before trusting absolute dollar figures, even a single aggregate observation, as
Section \ref{sec:depth-curve} does; do not assume that a dimensionally correct curve is
automatically a calibrated one. Fourth, design adaptation measures to cover distinct
CLIMADA mechanisms, since varying the mechanism, not just the numeric value of the
parameters, is what makes the analysis informative for the decision maker. Fifth, refine the
parametrization against the three acceptance criteria in Section \ref{sec:refinement}
before treating the final numbers as final, and resist the temptation to retune costs
purely to make an uncomfortable result look better.

\bibliographystyle{plainnat}
\bibliography{references}

\appendix

\section{Technical Notes and Recommendations for the Starter Pack}
\label{app:notes}

The following notes were compiled while extending the pack for this exercise, for
discussion with the UNU-EHS team and for other participants working on the Brazil case
study or extending their own. They cover friction points in the Starter Pack's data
structures, plotting utilities, and calibration workflow that were not obvious from the
existing teaching material, but that materially affected how long this exercise took and
how easy the results are to trust. Some are narrow, single-line fixes; others, such as the
distinction between a dimensionally correct curve and a locally validated one, point to a
broader methodological habit worth building into the pack's documentation rather than
left for each participant to rediscover independently. They are ordered roughly by how
early in a typical extension of the pack a participant is likely to encounter them.

\begin{enumerate}
\item \textbf{Exposures-swap measures interact badly with economic growth, and the pack
does not warn about it.} The \texttt{assets file} measure column points to a single HDF5
file, but \texttt{CostBenefit.calc} applies the present Entity's measures to the present
scenario and the future Entity's measures to the future scenario. If both Entities'
relocation measure point at the same, present-values HDF5 file, the measure's benefit
silently includes the removal of decades of economic growth from the adapted future
scenario, overstating it substantially. This exercise handled it by writing a second,
growth-scaled HDF5 file and repointing the future entity's measure in code. A short warning
in the cost-benefit notebook, or a worked example, would save other participants from this
trap.

\item \textbf{No worked example derives insurance parameters from the model itself.} The
training material introduces \texttt{Impact.calc\_freq\_curve()} and the risk transfer
attachment and cover parameters separately, but does not connect them. This exercise set
attachment and cover from the present-day unadapted exceedance curve, which felt like the
natural pattern to teach; it is worth adding to the notebook, since arbitrary first-guess
constants produce very different premiums and benefits.

\item \textbf{Relative paths in the \texttt{assets file} column resolve against the process
working directory, not the Entity file's own location.} This works while notebooks run from
the \texttt{notebooks/} folder, but breaks the moment an Entity file is loaded from
anywhere else. This friction recurred twice during this session's own corrections: once
when copying a corrected Entity file between two differently named folder layouts, and
again when renaming those same folders for public release, both times breaking the
relocation measure silently until re-executed and checked. Resolving \texttt{assets file}
relative to the Entity file's own location, rather than the process working directory,
would remove this entire class of bug.

\item \textbf{A dimensionally correct curve is not the same as a validated curve.} The pack
is careful about calibration uncertainty where it calibrates explicitly, but the
depth-based curve shipped with the Schools entity arrives pre-made, with no note of its
provenance or of whether its amplitude was ever checked against local data. Adopting the
equivalent literature curve for Residential and Companies and then checking it against
Porto Alegre's real 2024 loss (Section \ref{sec:depth-curve}) showed the unscaled version
overestimating that loss by roughly 20.7 times for Residential and 2.8 times for Companies.
Labelling the shipped curve with its source and adding a single sentence next to any
literature-curve suggestion, stating that it is a regional default whose amplitude must be
validated against available local loss data before absolute figures are trusted, would help
other participants catch this earlier.

\item \textbf{The expert-judgement bounds that resolve the calibration's degeneracy are
asset-specific, and nothing warns you when they stop applying.} The pack handles this
well: \texttt{06a} explains clearly that fitting two parameters to one observation leaves a
whole family of equally good solutions, and resolves it with stated expert judgement for
Schools, \texttt{max\_damage > 0.7} because several schools were rebuilt outright, and
\texttt{max\_intensity < 30} read off the parameter-space chart as the consequence of that
first choice. The difficulty is that the second bound is derived from the first, and only
the second survives transfer to another asset: Residential and Companies calibrate to
damage ratios of 0.028 and 0.294, so \texttt{max\_damage > 0.7} is impossible by
construction, while \texttt{max\_intensity < 30} still runs and still silently selects a
result. A participant adapting \texttt{06a} to a new asset therefore inherits a numerical
bound stripped of its justification. Restating the expert-judgement cell as a short
checklist, asking which bounds are asset-specific and what evidence would set them for the
new asset, would make that dependency visible at the point where it matters.

\item \textbf{\texttt{plot\_basemap\_eai\_exposure()}'s automatic colorbar collides with a
custom title once values exceed $10^6$.} Matplotlib adds a scientific-notation offset label
above the colorbar, at the same height as a centered axis title; for Schools, values under
200,000, this never appears, but Residential and Companies, values up to several million
USD, triggered it. The fix is to grab the colorbar axis and force a plain
\texttt{ScalarFormatter} instead of fighting it with title length or position. Also,
calling \texttt{ax.set\_title(..., loc="left")} after a function that already sets its own
centered title creates a second, independent title slot rather than replacing the first,
since matplotlib stores left, center, and right titles separately; call \texttt{set\_title}
without a \texttt{loc} argument to replace it.

\item \textbf{Windows-specific kernel failure, not reproduced on Linux.} On Windows, the
conda-forge \texttt{pyzmq}/\texttt{libzmq} build can make every
\texttt{jupyter nbconvert --execute} fail with \texttt{Bad file descriptor}
(\texttt{epoll.cpp:73}); replacing \texttt{pyzmq} with the PyPI wheel
(\texttt{pip install --force-reinstall pyzmq==26.2.0}) restores kernel communication. This
exercise's own notebooks were re-executed and verified on Linux without encountering this
issue, so it is specific to the Windows build of the conda-forge package, not to the
notebooks themselves; it is still worth a line in the installation FAQ for participants who
work on Windows.

\item \textbf{Insurance's benefit-to-cost ratio is approximately one over the cost factor
by construction, and that is worth stating explicitly.} With a premium loading of 1.5, the
ratio lands near 0.66 regardless of the layer. The measure's value is transferring
volatility, protection against rare, budget-breaking years, not expected-value savings; a
sentence in the teaching material would preempt an "insurance looks bad" misreading of the
bar chart.

\item \textbf{\texttt{06a}'s Excel export step is manual and mislabels the intensity
unit.} The final cell of \texttt{06a\_calibration\_schools.ipynb} prints the calibrated
functions with column \texttt{intensity\_m} and \texttt{intensity\_unit: "meter"}, even
though the calibration ran on flood duration in days, and the table must then be pasted
into the calibrated workbook by hand. Notebook \texttt{06b} automated this export with
openpyxl and wrote the unit as \texttt{days}. It is also worth clarifying that the
\emph{shipped} calibrated Schools entity uses a depth-based curve consistent with the JRC
return-period rasters, which is not the duration-based curve \texttt{06a} itself produces;
the connection between the two would be good to explain explicitly in the teaching
material.
\end{enumerate}

\end{document}
